\documentclass[10pt,a4paper]{amsart}
\usepackage{amsmath,amssymb,amsthm,graphicx}
% For commutative diagrams we use Xy-pic
\usepackage[all]{xy}

% We use 2cell for 2-commutative diagrams.
\xyoption{2cell}
% Theorem environments.
%
\theoremstyle{plain}
\newtheorem{theorem}[subsection]{Theorem}
\newtheorem{proposition}[subsection]{Proposition}
\newtheorem{lemma}[subsection]{Lemma}

\theoremstyle{definition}
\newtheorem{definition}[subsection]{Definition}
\newtheorem{example}[subsection]{Example}
\newtheorem{exercise}[subsection]{Exercise}
\newtheorem{situation}[subsection]{Situation}

\theoremstyle{remark}
\newtheorem{remark}[subsection]{Remark}
\newtheorem{remarks}[subsection]{Remarks}

\numberwithin{equation}{subsection}

% Macros
%
\def\lim{\mathop{\mathrm{lim}}\nolimits}
\def\colim{\mathop{\mathrm{colim}}\nolimits}
\def\Spec{\mathop{\mathrm{Spec}}}
\def\Hom{\mathop{\mathrm{Hom}}\nolimits}
\def\Ext{\mathop{\mathrm{Ext}}\nolimits}
\def\SheafHom{\mathop{\mathcal{H}\!\mathit{om}}\nolimits}
\def\SheafExt{\mathop{\mathcal{E}\!\mathit{xt}}\nolimits}
\def\Sch{\mathit{Sch}}
\def\Mor{\mathop{\mathrm{Mor}}\nolimits}
\def\Ob{\mathop{\mathrm{Ob}}\nolimits}
\def\Sh{\mathop{\mathit{Sh}}\nolimits}
\def\NL{\mathop{N\!L}\nolimits}
\def\CH{\mathop{\mathrm{CH}}\nolimits}
\def\proetale{{pro\text{-}\acute{e}tale}}
\def\etale{{\acute{e}tale}}
\def\QCoh{\mathit{QCoh}}
\def\Ker{\mathop{\mathrm{Ker}}}
\def\Im{\mathop{\mathrm{Im}}}
\def\Coker{\mathop{\mathrm{Coker}}}
\def\Coim{\mathop{\mathrm{Coim}}}

% Boxtimes
%
\DeclareMathSymbol{\boxtimes}{\mathbin}{AMSa}{"02}

%
% Macros for moduli stacks/spaces
%
\def\QCohstack{\mathcal{QC}\!\mathit{oh}}
\def\Cohstack{\mathcal{C}\!\mathit{oh}}
\def\Spacesstack{\mathcal{S}\!\mathit{paces}}
\def\Quotfunctor{\mathrm{Quot}}
\def\Hilbfunctor{\mathrm{Hilb}}
\def\Curvesstack{\mathcal{C}\!\mathit{urves}}
\def\Polarizedstack{\mathcal{P}\!\mathit{olarized}}
\def\Complexesstack{\mathcal{C}\!\mathit{omplexes}}
% \Pic is the operator that assigns to X its picard group, usage \Pic(X)
% \Picardstack_{X/B} denotes the Picard stack of X over B
% \Picardfunctor_{X/B} denotes the Picard functor of X over B
\def\Pic{\mathop{\mathrm{Pic}}\nolimits}
\def\Picardstack{\mathcal{P}\!\mathit{ic}}
\def\Picardfunctor{\mathrm{Pic}}
\def\Deformationcategory{\mathcal{D}\!\mathit{ef}}
\title{A simple root of a polynomial over a local ring is unique among the lifts congruent to it modulo the maximal ideal}
\author{The Stacks Project Authors}
\date{stacks-project:algebra.tex:local-snapshot-20261012; GNU FDL 1.2 or later}
\begin{document}
\maketitle
\noindent\textit{Editorial note.} All delivered mathematical text is the author's own, copied from the pinned Stacks Project chapter file algebra.tex; no statement or proof step was rewritten, reordered or removed.\par\medskip
\noindent\textit{Editorial note.} Preamble provenance, as required for this source: the first line of the pinned chapter file is \textbackslash{}input\{preamble\}, and the native preamble was not staged with the source, so no preamble text lives inside the chapter. The retained preamble used here is the project's single Stacks-wide preamble (the same 217-line / 8116-byte file that the already published chow and etale-cohomology cards use), copied mechanically line for line from sources/stacks-chow/preamble.tex lines 14-18 (Xy-pic setup) and 154-217 (theorem declarations and mathematical macros). Not one character was changed, no package was added, no macro was defined, and the excerpt serves only as preamble, never as a span coordinate.\par\medskip
\noindent\textit{Editorial note.} The same-chapter pointers are located at their \textbackslash{}label lines inside the pinned chapter file; the cross-chapter pointers are located in the project's own tag list staged in the snapshot (sources/stacks-chow/tags.txt), at the verbatim lines carrying their canonical tags.\par\medskip
\noindent\textit{Editorial note.} The author leaves the Taylor expansion to the reader ('details omitted'); that sentence is kept verbatim inside the counted span and is a routine omission, since the argument that uses the expansion is fully written out.\par\medskip
\noindent\textit{Editorial note.} The staged candidate list records ext\_refs = 0 for this label and the counted proof indeed contains no \textbackslash{}ref and no \textbackslash{}cite, so the card needs no navigation stub and no bibliography.\par\medskip
\noindent\textit{Allowed standard prerequisites (editorial).}\begin{itemize}
\item The Taylor expansion of a polynomial in two variables, f(x+y) = f(x) + f'(x)y + g(x,y)y\textasciicircum{}2. Used inside the counted proof at its first sentence. Supplies the identity that converts f(b) - f(a) into f'(a)(b-a) plus a multiple of (b-a)\textasciicircum{}2.
\item In a local ring an element of the form 1 + m with m in the maximal ideal is a unit. Used inside the counted proof in its last two sentences. Turns the factor 1 + f'(a)\textasciicircum{}\{-1\}c(b-a) into a unit so that the product vanishes only if b - a does.
\end{itemize}\medskip
\begin{lemma}
\label{lemma-uniqueness}
Let $(R, \mathfrak m, \kappa)$ be a local ring.
Let $f \in R[T]$. Let $a, b \in R$ such that $f(a) = f(b) = 0$,
$a = b \bmod \mathfrak m$, and $f'(a) \not \in \mathfrak m$.
Then $a = b$.
\end{lemma}

\begin{proof}
Write $f(x + y) - f(x) = f'(x)y + g(x, y) y^2$ in $R[x, y]$ (this is possible
as one sees by expanding $f(x + y)$; details omitted).
Then we see that $0 = f(b) - f(a) = f(a + (b - a)) - f(a) =
f'(a)(b - a) + c (b - a)^2$ for some $c \in R$. By assumption
$f'(a)$ is a unit in $R$. Hence $(b - a)(1 + f'(a)^{-1}c(b - a)) = 0$.
By assumption $b - a \in \mathfrak m$, hence $1 + f'(a)^{-1}c(b - a)$
is a unit in $R$. Hence $b - a = 0$ in $R$.
\end{proof}

\end{document}
